\section{GMM \& MLE}

\subsection{Gaussian Mixture Model (GMM)}

A \hl{Gaussian Mixture Model} assumes that each group has a multivariate
Gaussian distribution:

$$
\boxed{
\mathbf{X}\mid Z=k
\sim
\mathcal{N}(\boldsymbol{\mu}_k,\Sigma_k)
}
$$

Thus,

$$
f_k(\mathbf{x})
=
\mathcal{N}(\mathbf{x};
\boldsymbol{\mu}_k,\Sigma_k).
$$

The marginal mixture density is

$$
\boxed{
f(\mathbf{x})
=
\sum_{k=1}^K
p_k\mathcal{N}(\mathbf{x};
\boldsymbol{\mu}_k,\Sigma_k)
}
$$

Each group is characterized by

$$
\boxed{
p_k,\quad\boldsymbol{\mu}_k,\quad\Sigma_k
}
$$

where $p_k$ is the mixing proportion, $\boldsymbol{\mu}_k$ the mean, and
$\Sigma_k$ the covariance matrix.

\subsection{Maximum Likelihood Estimation (MLE)}

Given observations $\mathbf{x}_1,\ldots,\mathbf{x}_n$, MLE chooses the
parameters that make the observed data most likely:

$$
\boxed{
\hat{\theta}
=
\arg\max_\theta L(\theta)
}
$$

where the likelihood is

$$
L(\theta)
=
\prod_{i=1}^n f(\mathbf{x}_i;\theta).
$$

\subsection{Log-Likelihood}

Because products are difficult to optimize, we use the logarithm:

$$
\boxed{
\ell(\theta)
=
\log L(\theta)
=
\sum_{i=1}^n
\log f(\mathbf{x}_i;\theta)
}
$$

Since $\log$ is increasing,

$$
\arg\max_\theta L(\theta)
=
\arg\max_\theta\ell(\theta).
$$

\subsection{Empirical Estimates}

If the group memberships are known, let

$$
\tilde Z_{ik}=
\begin{cases}
1 & \text{if observation }i\text{ belongs to group }k,\\
0 & \text{otherwise.}
\end{cases}
$$

Then

$$
n_k=\sum_{i=1}^n\tilde Z_{ik}
$$

is the number of observations in group $k$.

The usual empirical estimates are

$$
\boxed{
\hat p_k=\frac{n_k}{n}
}
$$

$$
\boxed{
\hat{\boldsymbol{\mu}}_k
=
\frac{1}{n_k}
\sum_{i=1}^n
\tilde Z_{ik}\mathbf{x}_i
}
$$

and

$$
\boxed{
\hat\Sigma_k
=
\frac{1}{n_k}
\sum_{i=1}^n
\tilde Z_{ik}
(\mathbf{x}_i-\hat{\boldsymbol{\mu}}_k)
(\mathbf{x}_i-\hat{\boldsymbol{\mu}}_k)^T
}
$$

\textbf{Important:} These estimates assume that the group memberships
$\tilde Z_{ik}$ are known. In unsupervised GMM clustering, the memberships
are unknown and must be estimated, which leads to methods such as
Expectation-Maximization (EM).

\subsection{Key Connections}

$$
\boxed{
\text{GMM}
\rightarrow
p_k,f_k(\mathbf{x})
\rightarrow
P(Z=k\mid\mathbf{x})
\rightarrow
\text{classification}
\rightarrow
\text{risk}
\rightarrow
\text{Bayes optimal rule}
}
$$

$$
\boxed{
\text{MLE}
\rightarrow
\text{estimate }p_k,\boldsymbol{\mu}_k,\Sigma_k
\rightarrow
\text{fit the GMM}
}
$$

\textbf{Memory aids:}
\begin{itemize}
\item \textbf{Classification rule:} ``Which group do I assign $\mathbf{x}$ to?''
    \item \textbf{Risk:} ``How costly is my decision?''
\item \textbf{Bayes rule:} ``Choose the decision with the smallest expected cost.''
    \item \textbf{Equal costs:} ``Choose the group with the highest posterior probability.''
\item \textbf{GMM:} ``Each group is a Gaussian with its own $p_k,\mu_k,\Sigma_k$.''
    \item \textbf{MLE:} ``Choose parameters that make the observed data most likely.''
\end{itemize}